% =====================================================================
% ESQUELETO GENERAL DE LA TESIS — solo capítulos y secciones principales.
% Los subapartados de segundo nivel se añadirán capítulo a capítulo.
% El desglose detallado del Capítulo 4 vive en chapter_implementation.tex.
% =====================================================================


% =====================================================================
\chapter{Introduction}
% Se escribe al final. Fija el problema, las preguntas de investigación y
% las contribuciones que los capítulos 3, 4 y 6 deben sostener.
% =====================================================================

\section{Motivation}
\section{Problem Statement}
\section{Objectives and Research Questions}
\section{Contributions}
\section{Structure of the Thesis}


% =====================================================================
\chapter{Background and Related Work}
% Dos mitades: fundamentos necesarios para entender la propuesta, y estado
% del arte que justifica que la propuesta es necesaria. Cierra con el hueco
% de investigación que la tesis ocupa.
% =====================================================================

\section{Vehicular Telemetry: Data Sources and Service Requirements}
\section{Heterogeneous Vehicular Access: Cellular 5G NR and IEEE 802.11}
\section{Energy Consumption of Wireless Interfaces in Vehicular Nodes}
\section{Opportunistic Offloading and Connectivity Prediction}
\section{Deadline-Aware and Freshness-Constrained Data Delivery}
\section{Simulation Environments for Heterogeneous Vehicular Networks}
\section{Research Gap and Positioning of This Work}


% =====================================================================
\chapter{System Model and Decision Policy Design}
% Capítulo portante de la contribución teórica. Aquí viven todas las
% variables, ecuaciones y criterios que el Capítulo 4 solo realiza.
% =====================================================================

\section{Scenario and System Model}
\section{Assumptions and Design Constraints}
\section{Characterisation of WiFi Communication Opportunities}
\section{Prediction of the Time to the Next Communication Opportunity}
\section{Buffer Survival Condition and Minimum Offloading Volume}
\section{Deadline-Driven Rescue Condition}
\section{Complete Formulation of the Offloading Policy}


% =====================================================================
\chapter{Implementation}
% Estrictamente artefactos construidos. Desglose completo de subapartados
% en chapter_implementation.tex.
% =====================================================================

\section{Implementation Overview}
\section{Mobility Scenario Construction}
\section{Hybrid Vehicular Node}
\section{Event-Driven Energy Consumption Model}
\section{Hybrid Telemetry Application}
\section{Prediction Layer Implementation}
\section{Offloading Policy Implementation}
\section{Configuration Management and Experimental Control}
\section{Instrumentation and Data Collection}
\section{Verification of the Implementation}
\section{Implementation Constraints and Deviations from the Design}


% =====================================================================
\chapter{Evaluation Methodology}
% Define QUÉ se compara y CÓMO se mide. Ningún resultado numérico aquí.
% =====================================================================

\section{Evaluated Configurations and Baselines}
\section{Simulation Scenarios and Operating Conditions}
\section{Delivery and Reliability Metrics}
\section{Latency and Information Freshness Metrics}
\section{Energy Metrics}
\section{Experimental Design and Statistical Treatment}


% =====================================================================
\chapter{Results and Discussion}
% Una sección por eje de evaluación, y una final que los cruza. La comparación
% entre las cuatro configuraciones es el hilo conductor de todo el capítulo.
% =====================================================================

\section{Delivery Performance}
\section{Latency and Information Freshness}
\section{Buffer Occupancy and Data Loss}
\section{Energy Consumption}
\section{Energy--Freshness Trade-off Across Configurations}
\section{Discussion}
\section{Threats to Validity}


% =====================================================================
\chapter{Conclusions and Future Work}
% Responde una por una a las preguntas del Capítulo 1. Sin datos nuevos.
% =====================================================================

\section{Summary of Contributions}
\section{Answers to the Research Questions}
\section{Limitations of the Study}
\section{Future Work}
